%%% FIELD proof-packing-same-radius
Put \(m=n_1\) and \(\ell=\log(en)\).  If \(m=0\), the premise
\(mu\geq C\ell\) in \((\mathrm{P1})\) is impossible.  Moreover, the
definitions give \(q_t^L=0\) and \(q_t^U=1\), so \((\mathrm{P2})\) holds.
Thus it remains to consider \(m\geq1\).

For \(t\in\{0,1\}\), \(s\in[-1/2,1/2]\), and \(h\in H_n\), define
\[
 \mathcal C_{t,s,h}
 =\{z\in A_t:\|z-(0,s)\|_2\leq h\}.
 \tag{P4}
\]
For fixed points \(z_i=(r_i,s_i)\), membership of \(z_i\) in
\(\mathcal C_{t,s,h}\), viewed as a function of \(s\), is identically zero
unless \(z_i\in A_t\) and \(|r_i|\leq h\); otherwise it is the indicator of
the closed interval
\[
 [s_i-\sqrt{h^2-r_i^2},\ s_i+\sqrt{h^2-r_i^2}].
 \tag{P5}
\]
For a fixed pair \((t,h)\), the at most \(2N\) distinct endpoints in
\((\mathrm{P5})\) determine at most \(2N+1\) open cells and at most \(2N\)
endpoint traces.  The latter term is needed when left and right endpoints
tie.  Consequently the shatter coefficient of
\[
 \mathcal C_n=\{\mathcal C_{t,s,h}:t\in\{0,1\},
 s\in[-1/2,1/2],\ h\in H_n\}
\]
satisfies
\[
 \Pi_{\mathcal C_n}(N)
 \leq2|H_n|(4N+1)\leq C_1(N+1)\ell .
 \tag{P6}
\]

We record the relative empirical inequality used below and derive it here.
If \(P_m\) is the empirical law of \(m\) independent observations, then for
any measurable class \(\mathcal A\), \(x\geq1\), and
\(\Gamma=x+\log\{8\Pi_{\mathcal A}(2m)\}\),
\[
 \Pr\!\left\{\exists C\in\mathcal A:\ |P_m(C)-P(C)|>
 \sqrt{\frac{2P_m(C)\Gamma}{m}}+\frac{7\Gamma}{3m}\right\}
 \leq2e^{-x}.
 \tag{P7}
\]
To verify \((\mathrm{P7})\), introduce an independent empirical law
\(P'_m\), pair the two samples, and condition on the pooled pairs.  Randomly
swapping the two members of each pair represents
\(P_m(C)-P'_m(C)\) as a sum of independent centered signed variables, each
bounded by \(m^{-1}\), with conditional variance at most
\(2\bar p_C/m\), where
\(\bar p_C=\{P_m(C)+P'_m(C)\}/2\).  The elementary inequality
\[
 e^v\leq1+v+\frac{v^2}{2(1-|v|/3)},\qquad |v|<3,
\]
and independence imply, for \(0<\lambda<3m\),
\[
 \log E\!\left[
 e^{\lambda\{P_m(C)-P'_m(C)\}}\mid\text{pooled pairs}\right]
 \leq\frac{\lambda^2(2\bar p_C/m)}{2\{1-\lambda/(3m)\}}.
\]
Exponential Markov's inequality, optimization in \(\lambda\), and the same
calculation for the negative tail therefore give, for every pooled trace and
every \(z>0\),
\[
 \Pr\!\left\{|P_m(C)-P'_m(C)|>
 2\sqrt{\frac{\bar p_Cz}{m}}+\frac{2z}{3m}
 \mathrel{\Big|}\text{pooled pairs}\right\}\leq2e^{-z}.
 \tag{P8}
\]
There are at most \(\Pi_{\mathcal A}(2m)\) pooled traces, so a union bound
applies to \((\mathrm{P8})\).  We now eliminate the ghost count.  Apply the
same displayed moment bound to the unused half and separate the cases
\(P_m(C)\geq P(C)\) and \(P_m(C)<P(C)\).  In the first case substitute
\(P(C)\leq P_m(C)\); in the second substitute
\(P(C)\leq P_m(C)+d\), where \(d=|P_m(C)-P(C)|\).  In the latter case use
\[
 2\sqrt{d\Gamma/m}\leq d/2+2\Gamma/m.
\]
Solving the resulting two quadratic inequalities shows that failure of the
bound in \((\mathrm{P7})\) is contained in the union of at most four upper
or lower allocation-tail events of the form \((\mathrm{P8})\), with total
probability at most
\[
 8\Pi_{\mathcal A}(2m)e^{-\Gamma}=e^{-x}.
\]
This proves \((\mathrm{P7})\), with slack in its displayed factor two.  The
argument depends only on the finitely many pooled traces; increasing finite
trace subfamilies therefore also covers a nonattained supremum.

Apply \((\mathrm{P7})\) to \(\mathcal C_n\) with \(x=3\ell\).  By
\((\mathrm{P6})\), for all sufficiently large \(n\),
\(\Gamma\leq6\ell\).  Enlarging a numerical failure constant covers the
finitely many remaining \(n\), and therefore there is an event
\(\mathcal E_n\) such that, uniformly in the underlying law,
\[
 P(\mathcal E_n^c)\leq C_2n^{-3}
 \tag{P9}
\]
and, on \(\mathcal E_n\), simultaneously for every \(C\in\mathcal C_n\),
\[
 |P_m(C)-P(C)|\leq
 \sqrt{\frac{12P_m(C)\ell}{m}}+\frac{20\ell}{m}.
 \tag{P10}
\]
The closed intervals in \((\mathrm{P5})\) include every breakpoint trace,
including ties.  Since
\(P_m(\mathcal C_{t,s,h})=\widehat q_t^{(1)}((0,s),h)\), equation
\((\mathrm{P10})\) and the definitions of \(b_n,q_t^L,q_t^U\) imply
\[
 q_t^L(x,h)\leq q_t(x,h)\leq q_t^U(x,h)
 \tag{P11}
\]
simultaneously in \(x,t,h\).  Taking minima over \(t\) proves
\((\mathrm{P2})\).  If \(mq_t(x,h)\geq C_3\ell\), solving the two quadratic
inequalities in \((\mathrm{P10})\) gives
\[
 \tfrac12q_t(x,h)\leq q_t^L(x,h)\leq q_t(x,h)
 \leq q_t^U(x,h)\leq2q_t(x,h).
 \tag{P12}
\]
When \(ma_P(x,h)\geq C_3\ell\), both side masses satisfy the premise of
\((\mathrm{P12})\); taking minima proves \((\mathrm{P3})\).

Fix \(h\in H_n\).  Every endpoint in \((\mathrm{P5})\) generated by a
first-half score and lying in the target segment occurs in the source-labeled
list defining \(C_1(h)\).  Lexicographic ordering is total, and retaining
the least source label at a repeated coordinate retains that coordinate.
The same is true when a deterministic mesh point ties a breakpoint.  No
first-half membership indicator changes inside an open cell of the refined
list.  Since \(M_h=\lceil16/h\rceil\), consecutive deterministic mesh
points are separated by at most \(h/16\).  Thus the representative in the
definition satisfies, at cells and ties alike,
\[
 \|g_h(x)-x\|_2\leq h/32,
 \qquad
 \widehat q_t^{(1)}(g_h(x),h)=\widehat q_t^{(1)}(x,h).
 \tag{P13}
\]

Let \(F\) be an \(h\)-separated subset of
\(\{x\in B:a_P(x,h)\leq u\}\).  For each \(x\in F\), choose
\(t_x\) with \(q_{t_x}(x,h)\leq u\), which is possible by the definition
of \(a_P\).  Put \(z=\widehat q_{t_x}^{(1)}(x,h)\).  From
\((\mathrm{P10})\),
\[
 z\leq u+\sqrt{12z\ell/m}+20\ell/m
 \leq u+z/2+26\ell/m,
\]
and hence \(z\leq2u+52\ell/m\).  Substituting this bound into \(b_n(z)\),
using \(2\sqrt{uv}\leq u+v\), and invoking \((\mathrm{P13})\) yields
\[
 a^U(g_h(x),h)\leq C_4\{u+\ell/m\}.
 \tag{P14}
\]
For distinct \(x,x'\in F\), equation \((\mathrm{P13})\) gives
\[
 \|g_h(x)-g_h(x')\|_2
 \geq h-h/32-h/32=15h/16>h/2.
 \tag{P15}
\]
The upward dyadic rounding in the definition satisfies, for every
\(v\geq0\),
\[
 u(v)\leq\min\{1,2v+m^{-1}\},
 \qquad
 u(v)\geq\min\{1,v\}.
 \tag{P16}
\]
Choose the numerical constant \(c\geq C_4\).  If
\(c\{u+\ell/m\}\leq1\), equations \((\mathrm{P14})\) and
\((\mathrm{P16})\) make \(g_h(F)\) feasible in the empirical packing
problem defining \(\widehat J(h,c\{u+\ell/m\})\); if the argument exceeds
one, feasibility follows from \(a^U\leq1=u(c\{u+\ell/m\})\).
Equation \((\mathrm{P15})\) also makes the map injective.  Maximizing over
\(F\) proves
\[
 J_P(h,u)\leq\widehat J(h,c\{u+\ell/m\}).
 \tag{P17}
\]
This includes \(u=1\).

Conversely, let \(G=\{g_1<\cdots<g_M\}\) be a nonempty feasible set for
\(\widehat J(h,v)\), ordered along \(B\).  Equations \((\mathrm{P11})\)
and \((\mathrm{P16})\) imply
\[
 a_P(g,h)\leq a^U(g,h)\leq u(v)
 \leq\min\{1,2v+m^{-1}\},\qquad g\in G.
 \tag{P18}
\]
Feasibility gives consecutive gaps at least \(h/2\); collinearity of
\(B\) therefore gives
\[
 \|g_{j+2}-g_j\|_2
 =\|g_{j+2}-g_{j+1}\|_2+\|g_{j+1}-g_j\|_2\geq h.
 \tag{P19}
\]
The odd and even sublists are consequently \(h\)-separated at the same
physical radius \(h\), and \((\mathrm{P18})\) makes both feasible for the
population packing problem.  Hence
\[
 M\leq2J_P(h,\min\{1,2v+m^{-1}\}).
 \tag{P20}
\]
If the feasible maximum is empty, the defining floor of \(\widehat J\) is
one.  Substitute \(v=c\{u+\ell/m\}\) in \((\mathrm{P20})\), use monotonicity
of \(J_P(h,\cdot)\), and enlarge \(C\) to obtain
\[
 \widehat J(h,c\{u+\ell/m\})
 \leq C\{1\vee J_P(h,C\{u+\ell/m\})\}.
 \tag{P21}
\]
Together, \((\mathrm{P17})\) and \((\mathrm{P21})\) prove \((\mathrm{P1})\).
No comparison between different radii has been used.

Finally, the square-root endpoint map in the candidate list is continuous
on its closed domain.  Candidate inclusion, finite lexicographic sorting,
least-label duplicate removal, deterministic refinement, and midpoint
formation are Borel operations.  After absent candidate slots are encoded
by their Borel presence indicators, the finite maximum defining
\(\widehat J\) is a maximum of finitely many Borel feasibility indicators.
Thus \(C_1(h)\), \(G_1(h)\), and \(\widehat J(h,v)\) have the asserted
measurability, completing the proof.

%%% FIELD proof-profile-adaptive-estimator
Let \(\mathcal E_n\) be the first-half event supplied by
\(\mathrm{lem{:}empirical\text{-}packing\text{-}comparison\text{-}same\text{-}radius}\).
Uniformly over the declared union of base classes,
\[
 P(\mathcal E_n^c)\leq Cn^{-3}.
 \tag{T1}
\]
All lists and grids in the estimator are finite.  The measurability part of
that lemma makes the source-labeled sorting, refined grid, and representative
map Borel.  Dyadic severity rounding, the ordered greedy scan, the
first-covering-anchor rule, finite maximization over \(G_1(h)\), and
smallest-minimizer selection over the finite set \(H_n\) use finitely many
Borel comparisons and arithmetic operations.  Second-half counts and sums
are Borel, their zero-count branches are explicit, and both projections are
continuous.  Hence
\((O_1,\ldots,O_n,x)\mapsto\widehat\tau_n(x)\) is jointly Borel.  For each
sample, it is constant on finitely many open boundary cells with explicitly
assigned endpoints.  Since \(\tau_P\) is continuous on \(B\) by the
Lipschitz assumption, its supremum loss is the supremum over rational points
in those cells together with the finitely many endpoints, and is therefore
measurable.  The displayed zero-count and no-admissible-bandwidth clauses
make the estimator total.  Inspection of its construction shows that it
contains neither \(\kappa\) nor a profile-class label.

For \(h\in H_n\), define
\[
 z_P(x,h)=1+\log\!\left(
 e\{1\vee J_P(h,\min\{1,C_0a_P(x,h)\})\}\right)
\]
and
\[
 V_P(h)=\sup_{x\in B}\left[
 Lh+\sigma\sqrt{\frac{z_P(x,h)}{na_P(x,h)}}\right],
 \qquad
 \mathcal R_P=\min_{h\in H_n}V_P(h).
 \tag{T2}
\]
The base-mass assumption gives
\(a_P(x,h)\geq c_mh^\kappa>0\), while an \(h\)-separated subset of the
unit-length segment \(B\) is finite.  Thus every division and logarithm in
\((\mathrm{T2})\) is well defined, and the minimum is attained.  Choose
\(C_0\) no smaller than the numerical population-threshold constant in
\(\mathrm{lem{:}empirical\text{-}severity\text{-}net\text{-}extension}\).

We first compare criteria in the informative regime.  Suppose
\[
 n_1\inf_{x\in B}a_P(x,h)\geq C\log(en).
 \tag{T3}
\]
On \(\mathcal E_n\), equation \((\mathrm{P3})\) of the packing lemma, the
exact empirical-pattern identity at \(x\) and \(g_h(x)\), and the second
conclusion of the severity-net lemma imply, for every \(x\in B\),
\[
 a^L(g_h(x),h)\geq\tfrac12a_P(x,h),
 \qquad
 Z_n(g_h(x),h)\leq C z_P(x,h).
 \tag{T4}
\]
Here monotonicity of \(J_P(h,\cdot)\) justifies replacing the constant in
the severity-net lemma by \(C_0\).  Every \(g\in G_1(h)\) is its own
representative.  Therefore, if \((\mathrm{T3})\) holds and the resulting
lower bounds also imply admissibility, maximizing \((\mathrm{T4})\) over
\(g\in G_1(h)\) gives
\[
 \widehat V(h)\leq C V_P(h).
 \tag{T5}
\]

Let \(h_\star\) minimize \((\mathrm{T2})\), put
\(\ell=\log(en)\), and first suppose that
\[
 \mathcal R_P<c_\star\sigma/\sqrt\ell,
 \tag{T6}
\]
where \(c_\star>0\) is a sufficiently small numerical constant.  The
variance term at \(h_\star\) is at most \(\mathcal R_P\), and
\(z_P(x,h_\star)\geq1\), so
\[
 \frac{na_P(x,h_\star)}{z_P(x,h_\star)}
 \geq\frac{\sigma^2}{\mathcal R_P^2}
 >c_\star^{-2}\ell,
 \qquad x\in B.
 \tag{T7}
\]
For \(n\geq2\), both \(n_1\) and \(n_2\) are fixed positive fractions of
\(n\), up to numerical constants.  Taking \(c_\star\) sufficiently small,
\((\mathrm{T7})\) implies \((\mathrm{T3})\); then \((\mathrm{T4})\) gives
\[
 \frac{n_2a^L(g,h_\star)}{Z_n(g,h_\star)}
 \geq c\frac{na_P(g,h_\star)}{z_P(g,h_\star)}
 \geq2^{14},\qquad g\in G_1(h_\star).
\]
Thus \(h_\star\) is admissible, and the definition of the smallest
minimizer together with \((\mathrm{T5})\) yields
\[
 \widehat V(\widehat h)\leq\widehat V(h_\star)
 \leq C\mathcal R_P
 \quad\text{on }\mathcal E_n.
 \tag{T8}
\]

The complementary regime requires no additional law condition.  Suppose
\(\mathcal R_P\geq c_\star\sigma/\sqrt\ell\), and put
\(\rho_n=(\ell^2/n)^{1/\kappa}\).  For all sufficiently large \(n\), the
dyadic grid contains \(h_s\) such that
\[
 \rho_n\leq h_s<2\rho_n.
 \tag{T9}
\]
Existence is uniform for \(\kappa\in K\): the exponent window is compact,
\(\rho_n<h_0\) eventually, and \(\rho_n\) remains above the smallest grid
point \(h_02^{-\lceil\log_2n\rceil}\).  The base-mass assumption gives
\[
 \inf_{x\in B}a_P(x,h_s)\geq c_mh_s^\kappa
 \geq c_m\ell^2/n.
 \tag{T10}
\]
On \(\mathcal E_n\), equation \((\mathrm{P3})\) applies for all large
\(n\).  Independently of the law, an \(h_s/2\)-separated subset of the
unit-length segment has at most \(1+2/h_s\) points.  Hence the definition of
\(\widehat J\) gives
\[
 \max_{g\in G_1(h_s)}Z_n(g,h_s)
 \leq C\log(e/h_s)\leq C\ell.
 \tag{T11}
\]
Equations \((\mathrm{T10})\)--\((\mathrm{T11})\) imply
\(n_2a^L(g,h_s)\geq c\ell^2\geq2^{14}Z_n(g,h_s)\) for every grid point,
so \(h_s\) is admissible.  They also give
\[
 \widehat V(h_s)
 \leq CL\rho_n+C\sigma/\sqrt\ell
 \leq C\sigma/\sqrt\ell
 \leq C\mathcal R_P.
 \tag{T12}
\]
The penultimate inequality holds uniformly on
\(K\subseteq[2+\epsilon,\bar\kappa]\), because
\(n^{-1/\kappa}\ell^{2/\kappa+1/2}\to0\) uniformly on that compact window;
its constant may depend on the fixed ratio \(L/\sigma\).  Minimization again
gives \((\mathrm{T8})\).

Consequently, for all sufficiently large \(n\), on \(\mathcal E_n\) at
least one bandwidth is admissible.  On that event \(\widehat h\) is the
ordinary smallest minimizer, is measurable with respect to \(I_1\), and the
same selected bandwidth is used at every boundary point.

Invoke the conditional conclusion of
\(\mathrm{lem{:}empirical\text{-}severity\text{-}net\text{-}extension}\)
at this single \(I_1\)-measurable bandwidth.  The event \(\mathcal E_n\) is
also \(I_1\)-measurable, so the tower property and \((\mathrm{T8})\) yield
\[
 E_P\!\left[\|\widehat\tau_n-\tau_P\|_\infty
 \mathbf1_{\mathcal E_n}\right]
 \leq C\mathcal R_P.
 \tag{T13}
\]
On \(\mathcal E_n^c\), the trace-envelope assumption and the two
projections bound the loss by \(4L\).  Equations \((\mathrm{T1})\) and
\((\mathrm{T13})\), together with
\(n^{-3}\leq\log^2(en)/n\), therefore give
\[
 E_P\|\widehat\tau_n-\tau_P\|_\infty
 \leq C\mathcal R_P+C(L+\sigma)\frac{\log^2(en)}{n},
\]
which is precisely the asserted oracle inequality.  The minimization over
\(h\) remains outside the supremum over \(x\) because the construction uses
one global bandwidth.  For the finitely many sample sizes not covered by the
admissibility argument, both the estimator and the target lie in
\([-2L,2L]\), whereas \(\mathcal R_P\geq\sigma/\sqrt n\) because
\(z_P\geq1\) and \(a_P\leq1\).  Enlarging \(C\), uniformly over \(K\),
extends the same displayed inequality to every \(n\geq2\).  The explicit
fallback \(\widehat h=h_0\) handles any no-admissible-bandwidth sample and
is itself \(I_1\)-measurable.

%%% FIELD proof-endpoint-adaptation
The construction of \(\widehat\tau_n\) contains neither \(\kappa\) nor an
endpoint-class label.  We evaluate the oracle inequality from
\(\mathrm{thm{:}profile\text{-}adaptive\text{-}estimator}\) at one common
bandwidth for each class.

Fix \(\kappa\in K\).  For the isolated class, put
\(b=n^{-1/(\kappa+2)}\).  For all sufficiently large \(n\), uniformly over
the compact exponent window, choose \(h_b\in H_n\) with
\[
 b\leq h_b<2b.
 \tag{E1}
\]
For \(x\in B\), write
\[
 a_P(x,h_b)=w_xh_b^\kappa.
\]
The base-mass assumption gives \(w_x\geq c_m\).  Enlarge the fixed oracle
constant \(C_0\) so that \(C_0c_m\geq C_m\), and put
\(\alpha=\kappa-2\).  Eventually \(C_mh_b^\kappa\leq1\), uniformly on
\(K\).  If \(C_0w_xh_b^\kappa\leq1\), the isolated-envelope assumption,
applied at \(u=C_0a_P(x,h_b)\), gives
\[
 J_P(h_b,C_0a_P(x,h_b))
 \leq C_{\mathrm{iso}}\{1+Cw_x^{1/\alpha}\},
 \tag{E2}
\]
because
\[
 h_b^{-1}\left(\frac{C_0w_xh_b^\kappa}{h_b^2}\right)^{1/\alpha}
 =C_0^{1/\alpha}w_x^{1/\alpha}.
\]
If \(C_0w_xh_b^\kappa>1\), apply the same envelope at \(u=1\).  The case
premise implies
\[
 w_x^{1/\alpha}>C_0^{-1/\alpha}h_b^{-\kappa/\alpha}
 =C_0^{-1/\alpha}h_b^{-1-2/\alpha},
\]
so \((\mathrm{E2})\) remains valid when its left side is evaluated at the
clipped threshold \(\min\{1,C_0a_P(x,h_b)\}\).

Since \(\alpha\in[\epsilon,\bar\kappa-2]\), there is a finite constant,
uniform on the exponent window, such that
\[
 \sup_{\alpha\in[\epsilon,\bar\kappa-2]}
 \sup_{w\geq c_m}
 \frac{1+\log[e\{1+C(1+w^{1/\alpha})\}]}{w}<\infty.
 \tag{E3}
\]
Indeed, uniformly over this compact interval,
\(\log\{1+C(1+w^{1/\alpha})\}\leq C_K(1+\log_+w)\), and
\((1+\log_+w)/w\) is bounded on \([c_m,\infty)\).  Combining
\((\mathrm{E1})\)--\((\mathrm{E3})\), and using
\(nh_b^\kappa\asymp h_b^{-2}\), gives uniformly in \(x,P,\kappa\),
\[
 Lh_b+\sigma\sqrt{
 \frac{1+\log[e\{1\vee
 J_P(h_b,\min\{1,C_0a_P(x,h_b)\})\}]}
 {na_P(x,h_b)}}
 \leq Cb.
 \tag{E4}
\]
The theorem's remainder is \(o(b)\) uniformly on \(K\), because
\(\log^2(en)/n=o(n^{-1/(4+\epsilon)})\).  Taking the supremum over
\(x\), evaluating the theorem's bandwidth minimum at \(h_b\), and then
taking the supremum over the isolated class proves
\[
 \sup_{P\in\mathrm{Iso}_\kappa}
 E_P\|\widehat\tau_n-\tau_P\|_\infty
 \leq Cn^{-1/(\kappa+2)}.
 \tag{E5}
\]

For the pervasive endpoint, the required upper bound in fact holds on all
of \(\mathrm{Base}_\kappa\).  Every \(h\)-separated subset of the
unit-length segment \(B\) has cardinality at most
\(1+1/h\leq2/h\), since \(h\leq h_0\leq1/8\).  Therefore
\[
 1\vee J_P(h,u)\leq2/h
 \tag{E6}
\]
for every \(u\geq0\).  Put
\(r=(\log n/n)^{1/(\kappa+2)}\), and choose \(h_r\in H_n\) with
\(r\leq h_r<2r\), which is possible for all sufficiently large \(n\)
uniformly over \(K\).  The base-mass assumption and \((\mathrm{E6})\) give
\[
 \begin{aligned}
 &\sup_{x\in B}\left[
 Lh_r+\sigma\sqrt{
 \frac{1+\log[e\{1\vee
 J_P(h_r,\min\{1,C_0a_P(x,h_r)\})\}]}
 {na_P(x,h_r)}}\right]\\
 &\qquad\leq CLh_r+C\sigma
 \sqrt{\frac{\log(e/h_r)}{nc_mh_r^\kappa}}
 \leq C(\log n/n)^{1/(\kappa+2)}.
 \tag{E7}
 \end{aligned}
\]
Again the theorem remainder is of smaller order uniformly on \(K\).  Since
\(\mathrm{Perv}_\kappa\subseteq\mathrm{Base}_\kappa\), the theorem and
\((\mathrm{E7})\) yield
\[
 \sup_{P\in\mathrm{Perv}_\kappa}
 E_P\|\widehat\tau_n-\tau_P\|_\infty
 \leq C(\log n/n)^{1/(\kappa+2)}.
 \tag{E8}
\]
The identical total estimator occurs in \((\mathrm{E5})\) and
\((\mathrm{E8})\), so both endpoint bounds hold without preliminary regime
selection or knowledge of \(\kappa\).
